% !TeX program = xelatex
% !TeX program = xelatex
\documentclass[UTF8,a4paper]{ctexbook}
% The class determines paper geometry, type size, Chinese fonts, heading layout,
% double-sided pagination and chapter opening. No geometry, fancyhdr or font override.
\usepackage{amsmath,amssymb,amsthm,mathtools,mathrsfs}
\usepackage{tikz-cd}
\usepackage{enumitem,longtable,booktabs,array,needspace}
\usepackage[perpage,bottom]{footmisc}
\usepackage{graphicx,adjustbox}
\usepackage[unicode,hidelinks,bookmarksnumbered,hypertexnames=false]{hyperref}
\usepackage{bookmark}
% Retain EGA's historical chapter and paragraph numbering.
\ctexset{part={name={第,卷},number=\Roman{part}},chapter={number={\ifnum\value{chapter}=0 0\else\Roman{chapter}\fi}},section={name={\texorpdfstring{\S\,}{§},}}}
\renewcommand{\thesection}{\arabic{section}}
\renewcommand{\thesubsection}{\thesection.\arabic{subsection}}
\setcounter{secnumdepth}{2}
\setcounter{tocdepth}{2}
\setlength{\emergencystretch}{2em}
\allowdisplaybreaks[2]
\newcommand{\egapar}[1]{\par\addvspace{0.55\baselineskip}\Needspace{3\baselineskip}\noindent\hypertarget{para:\arabic{chapter}:#1}{}\hypertarget{entry:\arabic{chapter}:#1}{}\label{entry:\arabic{chapter}:#1}\textbf{（#1）}\quad\ignorespaces}
\newcommand{\egastatement}[2]{\par\addvspace{0.55\baselineskip}\Needspace{3\baselineskip}\noindent\hypertarget{para:\arabic{chapter}:#1}{}\hypertarget{entry:\arabic{chapter}:#1}{}\label{entry:\arabic{chapter}:#1}\textbf{#2（#1）}\quad\ignorespaces}
\newcommand{\eref}[1]{\hyperlink{entry:\arabic{chapter}:#1}{（#1）}}
\newcommand{\erritem}[2]{\par\addvspace{0.6\baselineskip}\Needspace{3\baselineskip}\noindent\hypertarget{errata:I:#1}{}\textbf{#2（#1）}\quad\ignorespaces}
\newcommand{\unnumberedchapter}[1]{\chapter*{#1}\addcontentsline{toc}{chapter}{#1}\markboth{#1}{#1}}
\newcommand{\src}[1]{\par\nobreak\noindent{\small 原文印刷页码：#1。}\par\nobreak\vspace{0.25\baselineskip}}
\newcommand{\origfootnote}[1]{\ifhmode\unskip\nobreak\fi\footnote{\textbf{原注：}#1}}
\newcommand{\transnote}[1]{\ifhmode\unskip\nobreak\fi\footnote{\textbf{译注：}#1}}
\newcommand{\separation}{\par\medskip\centerline{$*\qquad *\qquad *$}\medskip}
\newcommand{\F}{\mathscr F}
\newcommand{\G}{\mathscr G}
\newcommand{\HH}{\mathscr H}
\newcommand{\B}{\mathfrak B}
\newcommand{\OO}{\mathscr O}
\newcommand{\Kcat}{\boldsymbol K}
\newcommand{\restrict}[2]{#1\vert_{#2}}
\newcommand{\restmod}[2]{#1_{[#2]}}
\DeclareMathOperator{\Hom}{Hom}
\DeclareMathOperator{\Supp}{Supp}
\DeclareMathOperator{\Ker}{Ker}
\DeclareMathOperator{\Ann}{Ann}
\DeclareMathOperator{\Spec}{Spec}
\newcommand{\ilim}{\varinjlim}
\newcommand{\plim}{\varprojlim}
\newcommand{\A}{\mathscr A}
\newcommand{\BB}{\mathscr B}
\newcommand{\C}{\mathscr C}
\newcommand{\E}{\mathscr E}
\newcommand{\LL}{\mathscr L}
\newcommand{\M}{\mathscr M}
\newcommand{\N}{\mathscr N}
\newcommand{\PP}{\mathscr P}
\newcommand{\J}{\mathscr J}
\newcommand{\shHom}{\mathscr{H}\!om}
\newcommand{\dual}[1]{\check{#1}}
\DeclareMathOperator{\Coker}{Coker}
\let\Im\relax
\DeclareMathOperator{\Im}{Im}
\DeclareMathOperator{\Tor}{Tor}
\newcommand{\frJ}{\mathfrak J}
\newcommand{\frK}{\mathfrak K}
\newcommand{\frQ}{\mathfrak Q}
\newcommand{\frT}{\mathfrak T}
\newcommand{\frR}{\mathfrak R}
\newcommand{\frS}{\mathfrak S}
\newcommand{\comp}[1]{\widehat{#1}}
\newcommand{\ctensor}{\mathbin{\widehat{\otimes}}}
\DeclareMathOperator{\grad}{grad}
\newcommand{\wt}[1]{\widetilde{#1}}
\newcommand{\jj}{\mathfrak j}
\newcommand{\rr}{\mathfrak r}
\newcommand{\mm}{\mathfrak m}
\newcommand{\nrad}{\mathfrak N}
\newcommand{\aphi}{{}^{a}\varphi}
\DeclareMathOperator{\Spf}{Spf}
\newcommand{\frX}{\mathfrak X}
\newcommand{\frY}{\mathfrak Y}
\newcommand{\frZ}{\mathfrak Z}
\newcommand{\frU}{\mathfrak U}
\newcommand{\frV}{\mathfrak V}
\newcommand{\frW}{\mathfrak W}
\newcommand{\frD}{\mathfrak D}
\newcommand{\SSS}{\mathscr S}
\newcommand{\nn}{\mathfrak n}
\newcommand{\pp}{\mathfrak p}
\newcommand{\qq}{\mathfrak q}
\newcommand{\frj}{\mathfrak j}
\newcommand{\frk}{\mathfrak k}
\DeclareMathOperator{\Proj}{Proj}
\DeclareMathOperator{\Tr}{Tr}
\DeclareMathOperator{\Sym}{S}
\DeclareMathOperator{\gr}{gr}
\newcommand{\red}{\mathrm{red}}
\newcommand{\Gstar}{\Gamma_*}
\newcommand{\Tstar}{\Gamma_*}
\newcommand{\R}{\mathscr R}
\newcommand{\TT}{\mathscr T}
\newcommand{\I}{\mathscr I}
\newcommand{\K}{\mathscr K}
\newcommand{\V}{\mathbf V}
\newcommand{\Pb}{\mathbf P}
\newcommand{\Z}{\mathbf Z}
\newcommand{\Gr}{\operatorname{Gr}}
\newcommand{\id}{\mathrm{id}}
\DeclareMathOperator{\codim}{codim}
\DeclareMathOperator{\rk}{rg}
\DeclareMathOperator{\Pic}{Pic}
\newcommand{\sProj}{\mathscr{P}\!roj}
\newcommand{\catHom}{\mathscr{H}\!om}
\newcommand{\Ens}{\mathbf{Ens}}
\newcommand{\OI}{0_{\mathrm I}}
\newcommand{\OIII}{0_{\mathrm{III}}}
\newcommand{\Ext}{\operatorname{Ext}}
\newcommand{\shExt}{\mathscr{E}\!xt}
\newcommand{\shTor}{\mathscr{T}\!or}
\newcommand{\Hh}{\mathbb H}
\newcommand{\len}{\operatorname{long}}
\newcommand{\cplx}[1]{#1^{\bullet}}
\newcommand{\Cat}{\mathbf C}
\newcommand{\ra}{\longrightarrow}
\newcommand{\la}{\longleftarrow}
\newcommand{\iso}{\xrightarrow{\sim}}
\newcommand{\eproof}{\par\noindent\textbf{证明。}\ignorespaces}
\newcommand{\finishproof}{\hfill$\square$\par}
\newcommand{\srcpage}[1]{\par\noindent{\small\itshape〔底本第#1页〕}\par\nobreak}
\newcommand{\cocomp}[1]{\complement #1}
\DeclareMathOperator{\Aut}{Aut}
\DeclareMathOperator{\GL}{GL}
\DeclareMathOperator{\trdeg}{trdeg}

% Macros for the reviewed continuation. Existing mathematical macros are unchanged.
\newcommand{\kk}{\mathfrak k}
\newcommand{\cH}{\mathscr H}
\newcommand{\hTor}{\mathbb{T}\!or}
\newcommand{\lTor}{\mathscr{T}\!or}
\newcommand{\spE}{\mathscr E}
\DeclareMathOperator{\coprof}{coprof}
\newcommand{\source}[1]{\par\smallskip\noindent\textit{底本位置：#1。}\par\smallskip}
\newcommand{\newscope}{new}
\newcommand{\egap}[1]{\par\addvspace{.6\baselineskip}\Needspace{3\baselineskip}\noindent\hypertarget{new:\newscope:#1}{}\textbf{（#1）}\quad\ignorespaces}
\newcommand{\egas}[2]{\par\addvspace{.6\baselineskip}\Needspace{3\baselineskip}\noindent\hypertarget{new:\newscope:#1}{}\textbf{#2（#1）}\quad\ignorespaces}
\newcommand{\cmt}[2]{\footnote{\textbf{译校注〔#1〕：}#2}}
\newcommand{\ednote}[1]{\footnote{\textbf{英译编者注：}#1}}
\newcommand{\orig}[1]{\footnote{\textbf{原注：}#1}}
\newcommand{\head}[1]{\chapter*{#1}\addcontentsline{toc}{chapter}{#1}\markboth{#1}{#1}}
\newcommand{\egasection}[1]{\section*{#1}\addcontentsline{toc}{section}{#1}}
\newcommand{\subsec}[1]{\subsection*{#1}\addcontentsline{toc}{subsection}{#1}}
\newcommand{\audit}[3]{\Needspace{5\baselineskip}\hypertarget{audit:#1}{}\section*{#1\quad #2}\addcontentsline{toc}{section}{#1\quad #2}\noindent\textbf{判定：#3。}\par\smallskip}
\pdfstringdefDisableCommands{\def\quad{ }\def\qquad{ }\def\S{§}}

\newcommand{\bTor}{\operatorname{\mathbf{Tor}}}
\newcommand{\termTor}{\mathop{\mathscr{T}\!or^{\mathrm{term}}}\nolimits}
\newcommand{\NewNote}[1]{\par\smallskip\noindent\textbf{编排说明。}#1\par\smallskip}
\hypersetup{pdftitle={代数几何基础：按内容顺序编排的中文译校合订本（第七版）},pdfauthor={原著 A. Grothendieck；与 J. Dieudonné 合作；预备笔记英译编者见各篇},pdfsubject={第0章归并；正文连续；已译范围和未译范围分开列明}}
\title{代数几何基础\\中文译校合订本\\[1em]\large 按内容顺序编排·第七版}
\author{A. Grothendieck\\与 J. Dieudonné 合作撰写\\[.6em]\normalsize 预备笔记英译及编辑：Piotr Blass、Joseph Blass}
\date{}

% Additions for the fourth continuation; no font or page-geometry overrides.
\DeclareMathOperator{\Ass}{Ass}
\DeclareMathOperator{\prof}{prof}
\DeclareMathOperator{\htp}{ht}
\newcommand{\pd}{\operatorname{dim.proj}}
\newcommand{\injdim}{\operatorname{dim.inj}}
\newcommand{\gldim}{\operatorname{dim.coh}}
\newcommand{\Salg}{\mathrm S}

\DeclareMathOperator{\End}{End}

% Additions for chapter 0, section 18. Keep class defaults.
\DeclareMathOperator{\Exan}{Exan}
\DeclareMathOperator{\Exal}{Exal}
\DeclareMathOperator{\Exalcom}{Exalcom}
\DeclareMathOperator{\Exantop}{Exantop}
\DeclareMathOperator{\Exaltop}{Exaltop}
\DeclareMathOperator{\Exalcotop}{Exalcotop}
\newcommand{\HomCont}{\operatorname{Hom.cont}}

\begin{document}
\frontmatter
\maketitle
\head{本版编排与收录说明}
\noindent\textbf{这是一份按内容顺序整理的中文工作译稿，不是全部已上传材料的完结全译本。}第I、II章及第III章的两个已出版分册正文已经连续收入；第0章当前收入§§1--18及§19的19.0--19.6；第IV章主体、第0章19.7--19.10及§§20--23尚未连续译入；第V卷英文预备笔记只收入已译出的部分。未译内容没有用原文扫描页、提要或空白标题代替译文。

\section*{正文次序}
本版不再按上传时间、译稿批次或独立PDF页序排列，而按原著内容编号组织：导言；第0章预备知识；第I章概形的语言；第II章若干类态射的初等整体研究；第III章凝聚层的上同调研究；后续待译范围说明；独立编排的第V卷预备笔记已译部分；最后为各原分册的书末资料。第III章第一部分的§§1--5与第二部分的§§6--7直接衔接，不再被第一部分的索引隔开。

第0章原来分载于第一卷、第三卷第一部分、第四卷第一部分。本版将已译的§§1--7、§§8--13、§§14--18和§19已译部分归并在一起，但保留 $0_{\mathrm I}$、$0_{\mathrm{III}}$、$0_{\mathrm{IV}}$ 等出处记号。它们标示原出版分册，并非重复的三章。这样的安排依内容编号，而不声称每一节的全部前提均已在前文出现；原著跨章引用仍按原号保留。

\section*{实际收录范围}
\begin{longtable}{@{}p{.25\linewidth}p{.67\linewidth}@{}}
\toprule 材料 & 已收入的连续译文及状态\\\midrule\endhead
第一卷（1960） & 导言；第0章§§1--7；第I章§§1--10，止于（10.15.7）；原书书末资料。承接既有V2校订源码。\\
第二卷（1961） & 第II章§§1--8，止于（8.14.14）；原书参考文献、索引、目录、勘误清单一。其书末资料沿用上一合订版依据底本恢复的译排来源。\\
第三卷第一部分（1961） & 第0章§§8--13；第III章§§1--5，止于（5.5.2）；原书书末资料。\\
第三卷第二部分（1963） & 第III章§§6--7，止于（7.9.14），连同其证明、记号索引、术语索引、原目录及勘误清单二。第三版已补齐（6.7.1）至本分册末尾；本版保留。\\
第四卷第一部分（1964） & 第0章补篇章首说明、§§14--18及§19的19.0--19.6，止于（19.6.7）末尾。本版新增（19.5.1）--（19.6.7）；19.7以后尚未译入。\\
第V卷英文预备笔记 & \texttt{EGAV1-2.ps} 的六页可读材料；\texttt{EGAVaiello.pdf} 的卷首说明及“超平面截面与锥投影”§§1--6。余下内容不计入已译范围。\\
\bottomrule
\end{longtable}

\section*{译文、勘误与核查的关系}
正文采用本轮已经核实的正确陈述、公式和证明；原文错误写法、修改理由及依据记入脚注。作者正式勘误与本次译者补证分开，较长反例与验证收入独立核查报告。尚未确认的疑点不擅自改成确定结论；原著本来省略的证明、预备笔记中未写完或无法辨认的内容，不凭通常知识补造成“原文”。本轮新增注号以G7开头，旧注号保留以便追踪。

这一编辑原则已落实在本轮新增的19.5--19.6，并回溯修改了（19.3.5，v）、（19.3.10.1）、（19.3.12）和（19.4.1）。\textbf{它不表示此前全部旧译文中的历史核查意见均已逐项回写正文。}其他既有部分继续保存原稿；其校订状态须参看各条脚注和历史报告，后续仍需逐项核定。

对作者已撤销的（19.3.12），正文仅保留编号和撤销说明，原陈述及问题写入脚注；完整历史陈述与原证明保存在报告附录，不再占据有效论述的正文位置。作者勘误清单二、清单三相关选译仍在书末收录。未经核实的替代结论没有补造。

第V卷上传材料的标题明确标为预备笔记的翻译与编辑，英译编者为Piotr Blass与Joseph Blass。它们包含重组建议、未完成句子及待补证明，不按正式刊行的第I--IV章同等状态标示。与英文合并稿重合的PS载体不重复计作新译文。

\section*{第七版新增范围与回溯修订}
本版从第六版实际接续点（19.5.1）起，译入19.5“形式光滑性与相伴分次环”以及19.6“域上的代数的情形”，止于（19.6.7）末尾。对应第四卷第一部分分册印刷第90--104页（上传PDF第89--103页）。本轮包括一般定理的全部分段证明、Cohen定理、系数域和几何正则性判据；19.7章首及（19.7.1）仍为下一接续点，未算作译完。

（19.6.2）在正文补明相容的子域嵌入 $k\subset A$；原文仅给 $k\subset K$ 的问题及 $\mathbf Z_p$ 反例放在脚注和报告中。（19.5.6.7）的核比较步骤换成有限滤过逐级消去论证；（19.6.6）以Galois作用在局部分量上的传递性及忠实平坦下降修补原证明。后两项是证明修补，不宣称原命题结论有反例。

第0章其余部分、第I--III章及预备笔记的旧正文，除上列回溯项外没有进行一轮新的逐句复审。本次新核查范围只包括新增内容及上述直接复核的相关条目。累计核查报告中“第六版”“第五版”等章节是历史档案，其中旧的正文状态须以本版实际源码为准。

\section*{版式与使用}
采用 \texttt{\string\documentclass[UTF8,a4paper]\{ctexbook\}}，由文档类决定默认中文字体、字号、版心、标题与双面分页。不另设页边距，不用自定义页眉。为保留原著结构，仅保留原章号、节号与公式号；默认双面排版所需的空白偶数页不删除。前置部分使用罗马数字页码，其后连续编号；总目录指向本合订本页面，各原著目录的页码仍指原版。

PDF可用于对照研读；可独立编译的单文件TeX与分文件工程同时提供。旧稿内容的编排归并与技术校验，不等于本次将旧稿的每一句译文或每一个证明重新独立审定。累计核查报告单独交付，历史核查意见带有批次说明；对同一事项，本版直接核得的作者勘误和明确论证须与历史意见对照阅读，不能将历史“待核”标签当作最新状态。
\tableofcontents

\mainmatter
\input{text/introduction.tex}
\input{text/ch0_01_07.tex}
\input{text/ch0_08_13.tex}
\renewcommand{\newscope}{0-IV}
\input{text/ch0_14_15a.tex}
\input{text/ch0_15b.tex}
\input{continuations/iv0_16_12.tex}
\input{continuations/iv0_16_3.tex}
\input{continuations/iv0_16_45.tex}
\input{continuations/iv0_17_1.tex}
\input{continuations/iv0_17_2.tex}
\input{continuations/iv0_17_3.tex}
\input{continuations/iv0_18_1.tex}
\input{continuations/iv0_18_2.tex}
\input{continuations/iv0_18_3.tex}
\input{continuations/iv0_18_4.tex}
\input{continuations/iv0_18_5.tex}
\input{continuations/iv0_19_01.tex}
\input{continuations/iv0_19_23.tex}
\input{continuations/iv0_19_4.tex}
\input{continuations/iv0_19_5a.tex}
\input{continuations/iv0_19_5b.tex}
\input{continuations/iv0_19_6.tex}
\NewNote{本章本版连续译至§19的（19.6.7）末尾。19.7--19.10、第0章§§20--23尚未译入；下一接续点为19.7“局部同态的情形：存在性与唯一性定理”的章首及（19.7.1）。下接第I章不表示第0章余文已译完。}
\input{text/ch1.tex}
\input{text/ch2.tex}
\input{text/ch3_01_05.tex}
\renewcommand{\newscope}{III-2}
\input{text/ch3_06a.tex}
\input{text/ch3_06b.tex}
\input{text/ch3_06c.tex}
\input{continuations/iii2_67.tex}
\input{continuations/iii2_68.tex}
\input{continuations/iii2_69.tex}
\input{continuations/iii2_610.tex}
\input{continuations/iii2_71_72.tex}
\input{continuations/iii2_73.tex}
\input{continuations/iii2_74.tex}
\input{continuations/iii2_75_76.tex}
\input{continuations/iii2_77.tex}
\input{continuations/iii2_78_79.tex}
\head{第IV章及后续材料的待译范围}
\noindent\textbf{本页是范围说明，不是正文译文。}按原著内容顺序，第III章之后应接第IV章《概形与概形态射的局部研究》。当前已译的第四卷材料只是其所附第0章§§14--18及§19的19.0--19.6，已移到本书第0章；第IV章主体并未因标题或目录出现而被计作译完。

核对上传分册的原目录，剩余连续译文包括：第四卷第一部分的第0章19.7--19.10及§§20--23、第四章章首及§1；第二部分的第四章§§2--7；第三部分的§§8--15；第四部分的§§16--21及书末资料。各节内部仍须以原段落为单位翻译，不用提纲替代。第0章下一接续条目为19.7小节章首及（19.7.1）。

第V卷预备笔记不补作一个并不存在于当前译稿中的完整“第V章”。下编仅收录已译出的短篇及“超平面截面与锥投影”开头；该篇下一接续点为（7.1）。短篇中原有未完成句子仍按原状标示。

\part*{预备笔记：已译部分}
\addcontentsline{toc}{part}{预备笔记：已译部分}
\renewcommand{\newscope}{V-S}
\input{text/v1_2.tex}
\renewcommand{\newscope}{V-H}
\input{text/v5_01_02.tex}
\input{text/v5_03_06.tex}
\backmatter
\part*{各原分册的书末资料}
\addcontentsline{toc}{part}{各原分册的书末资料}
\input{backmatter/volume_I.tex}
\input{backmatter/volume_II.tex}
\input{backmatter/volume_III_1.tex}
\input{backmatter/volume_III_2.tex}
\input{backmatter/errata_III_2.tex}
\input{backmatter/bibliography_IV_1.tex}
\head{第四卷勘误清单三：与第0章§§16--19已译部分有关的选译}
\source{EGA IV第四部分，勘误清单三，印刷第345、347页；上传PDF第344、346页}
\NewNote{本节只收入与已译范围直接相关的七项作者勘误，保留其原分类和原编号。它不是整份勘误清单三的全译；未列出的条目不因此被删除或判定无关。下文的“页”均指作者所列原出版分册的印刷页码。}

\egasection{A）排印错误：三项选译}
\hypertarget{ErrIV:A16222}{}\noindent\textbf{（$0_{\mathrm{IV}}$，16.2.2.2）。}
第26页第13行，将 $M'\cap\mathfrak q^nM'$ 改为 $M'\cap\mathfrak q^nM$。

\medskip
\hypertarget{ErrIV:A1737}{}\noindent\textbf{（$0_{\mathrm{IV}}$，17.3.7）。}
第50页倒数第16行，将（17.1.10）改为（17.3.3）。

\medskip
\hypertarget{ErrIV:A1813}{}\noindent\textbf{（$0_{\mathrm{IV}}$，18.1.3）。}
在图（18.1.3.1）中，$p_2$ 应标在箭头 $G\to F$ 的旁边。

\egasection{B）正文修订：四项选译}
\hypertarget{ErrIV:5}{}
\noindent\textbf{（Err$_{\mathrm{IV}}$，5）。}
在（$0_{\mathrm{IV}}$，16.1.4）中，第23页倒数第8行，将“说 $A$ 是悬链的”改为：“说一个整环 $A$ 是悬链的”。

\medskip
\hypertarget{ErrIV:6}{}
\noindent\textbf{（Err$_{\mathrm{IV}}$，6）。}
在（$0_{\mathrm{IV}}$，16.1.9）中，可删去 $A$ 为诺特环的假设。事实上，对于每个 $u\in B_M$，$A[u]$ 是 $C$ 的一个子 $A$-代数，故 $u$ 整于 $A$（Bourbaki，\textit{交换代数}，第五章，§1，第1目，命题1的推论1），因而也整于 $A_M$；仍由（16.1.5）得出结论。

\medskip
\hypertarget{ErrIV:7}{}
\noindent\textbf{（Err$_{\mathrm{IV}}$，7）。}
在（$0_{\mathrm{IV}}$，16.5.12）中，应注意：关系（16.5.11.1）等价于将 $A$ 换成各商环 $A/\mathfrak r$ 后得到的相同关系，其中 $\mathfrak r$ 遍历 $A$ 的极小素理想。因此，在（16.5.12）的证明中可限于 $A$ 为整环的情形，并应用（16.1.4.2）。


\medskip
\hypertarget{ErrIV:8}{}
\noindent\textbf{（Err$_{\mathrm{IV}}$，8）。}
删去（$0_{\mathrm{IV}}$，19.3.12），因为其证明错误（集合 $W_\lambda,U_\lambda,V_\lambda$ 不是线性流形）。

\part*{译稿来源与版本记录}
\addcontentsline{toc}{part}{译稿来源与版本记录}
\markboth{译稿来源与版本记录}{译稿来源与版本记录}
\NewNote{下面三份说明来自既有分卷，保留其中关于当时稿件范围和制作过程的历史记录。旧页数、旧接续点不代表本合订本的当前页数或当前完成度；本版范围以卷首说明及工程包PROGRESS.md为准。}
\input{backmatter/old_note_I.tex}
\input{backmatter/old_note_II.tex}
\input{backmatter/old_note_III_1.tex}
\input{backmatter/additional_sources.tex}
\end{document}
